\chapter{All-Pairs Shortest Paths}

The previous chapters answered a single-source question: fix a vertex $s$ and find the
distance to every other vertex. Here we raise the stakes. We want the distance between
\emph{every} ordered pair $(u,v)$ at once -- a full $n \times n$ table of shortest-path
costs. This is the all-pairs shortest paths (APSP) problem, and it sits at a crossroads of
dynamic programming, matrix algebra, and clever reweighting tricks.

There are two ways to attack it. The first treats APSP as one big dynamic program over the
whole graph and produces the Floyd--Warshall algorithm in three nested loops -- fifteen lines
that handle negative edges, detect negative cycles, and double as a transitive-closure machine.
The second runs a fast single-source solver from every vertex, which is great on sparse graphs
but stumbles when edges go negative. Johnson's algorithm rescues the second approach by
\emph{reweighting} the graph so that all edges become nonnegative without changing which paths
are shortest -- a potential trick borrowed straight from the Bellman--Ford world. We will build
both, prove both, trace both on a tiny graph, and tie Floyd--Warshall to the min-plus matrix
product that connects shortest paths to linear algebra.

\chapterline

\section{The problem, and why it is not just $n$ Dijkstras}
\label{ch09:problem}

\begin{focusbox}
We are given a directed graph $\Gph=(V,E)$ with $n=\abs{V}$ vertices, $m=\abs{E}$ edges, and a
weight $w(u,v)\in\R$ on each edge. We want $\delta(u,v)$, the weight of a minimum-weight path
from $u$ to $v$, for all $n^2$ ordered pairs. Edge weights may be negative; if a negative-weight
cycle lies on a path between some pair, the relevant distances are $-\infty$ and we want to
\emph{report} that rather than return garbage.
\end{focusbox}

The naive reduction is tempting: run a single-source algorithm $n$ times, once per source. On a
nonnegative graph each Dijkstra run costs $\Oh(m + n\log n)$ with a Fibonacci heap, so the total
is $\Oh(nm + n^2\log n)$. For sparse graphs ($m=\Oh(n)$) that is $\Oh(n^2\log n)$, which is
excellent. But two things break this plan.

\begin{itemize}
  \item \textbf{Negative edges.} Dijkstra is wrong the moment an edge weight is negative, because
        its greedy ``settle the closest vertex'' invariant assumes distances only grow. We would
        fall back to Bellman--Ford at $\Oh(nm)$ per source, giving $\Oh(n^2 m)$ overall -- on a dense
        graph $m=\Th(n^2)$ that is $\Oh(n^4)$. Painful.
  \item \textbf{Dense graphs.} When $m=\Th(n^2)$, even $n$ Dijkstras cost $\Oh(n^3 + n^2\log n)=\Oh(n^3)$.
        At that density a single $\Oh(n^3)$ dynamic program with a tiny constant and contiguous memory
        access -- Floyd--Warshall -- is simpler and usually faster in practice.
\end{itemize}

\begin{intuitionbox}
Two regimes, two champions. \textbf{Dense graphs, or negative edges with simple code:}
Floyd--Warshall, $\Oh(n^3)$, no heaps, cache-friendly. \textbf{Sparse graphs:} reweight once with
Johnson's potentials, then run $n$ Dijkstras for $\Oh(nm + n^2\log n)$. The whole chapter is about
earning the right to make that choice.
\end{intuitionbox}

\begin{historybox}
Floyd--Warshall is a convergence of ideas. Bernard Roy (1959) and Stephen Warshall (1962) gave the
transitive-closure version; Robert Floyd (1962) adapted Warshall's idea to shortest paths in a
half-page \emph{Communications of the ACM} note. The shared loop structure was so natural that the
same recurrence was rediscovered repeatedly. Donald Johnson's reweighting algorithm came in 1977,
and it is the textbook example of using a feasible \emph{potential} to convert a hard-looking problem
into an easy one -- the same move that powers min-cost flow.
\end{historybox}

\section{Floyd--Warshall: shortest paths as a dynamic program}
\label{ch09:fw}

\subsection{The right subproblem}

The art of dynamic programming is choosing what to vary. For a single path the obvious knob is its
\emph{length} (number of edges), and indeed that gives a working but slower recurrence we meet in
Section~\ref{ch09:minplus}. Floyd and Warshall found a cleverer knob: restrict which
\emph{intermediate} vertices a path is allowed to use.

Number the vertices $1,2,\dots,n$. Define
\[
  d^{(k)}(i,j) = \text{weight of the shortest path from } i \text{ to } j
  \text{ whose intermediate vertices all lie in } \set{1,\dots,k}.
\]
``Intermediate'' means strictly between the endpoints; $i$ and $j$ themselves are always allowed.
When $k=0$ no intermediate vertices are permitted, so the only legal paths are single edges:
\[
  d^{(0)}(i,j) = \begin{cases}
    0 & i=j,\\
    w(i,j) & (i,j)\in E,\\
    +\infty & \text{otherwise.}
  \end{cases}
\]
When $k=n$ every vertex is allowed as an intermediate, so $d^{(n)}(i,j)=\delta(i,j)$ is exactly the
answer we want. The whole algorithm is the climb from $k=0$ to $k=n$.

\begin{intuitionbox}
Think of the vertices $1,\dots,n$ as doors that open one at a time. At stage $k$ the first $k$ doors
are open. A path may walk through any open door as a waypoint but may not yet use doors
$k{+}1,\dots,n$. Opening door $k$ can create only one genuinely new kind of shortcut: paths that
detour \emph{through} $k$. That single observation is the entire recurrence.
\end{intuitionbox}

\subsection{The recurrence}

Fix a pair $(i,j)$ and suppose we know all of $d^{(k-1)}$. Now we open door $k$. Consider any
shortest path $P$ from $i$ to $j$ using intermediates in $\set{1,\dots,k}$. Exactly one of two things
is true:

\begin{itemize}
  \item $P$ does not use vertex $k$ as an intermediate. Then $P$ already lives in $\set{1,\dots,k-1}$,
        so its cost is $d^{(k-1)}(i,j)$.
  \item $P$ uses vertex $k$ once as an intermediate. (Why only once? See the trap box.) Split $P$ at
        that occurrence into $i \rightsquigarrow k$ then $k \rightsquigarrow j$. Both halves use
        intermediates only in $\set{1,\dots,k-1}$, so their best costs are $d^{(k-1)}(i,k)$ and
        $d^{(k-1)}(k,j)$.
\end{itemize}

Taking the cheaper of the two options:
\begin{equation}
  d^{(k)}(i,j) = \min\!\Big(\, d^{(k-1)}(i,j),\;\; d^{(k-1)}(i,k) + d^{(k-1)}(k,j)\,\Big).
  \label{ch09:fw-recurrence}
\end{equation}

\begin{trapbox}
Why may we assume the through-$k$ path uses $k$ \emph{exactly once}? Because we can take $P$ to be a
shortest \emph{walk}, and if there are no negative cycles a shortest walk never repeats a vertex --
repeating $k$ would mean a cycle through $k$, which we could delete to get a no-heavier path. If a
negative cycle \emph{is} present the recurrence still runs, but a diagonal entry $d^{(k)}(i,i)$ goes
negative; detecting that is exactly how Floyd--Warshall flags negative cycles
(Section~\ref{ch09:fw-negcycle}). The crucial structural point: both halves $i\to k$ and $k\to j$ avoid
$k$ as an \emph{interior} vertex, so they are legitimately scored by $d^{(k-1)}$.
\end{trapbox}

\subsection{Collapsing the third dimension}

Written as stated, $d^{(k)}$ is a stack of $n+1$ matrices -- $\Th(n^3)$ memory. We can do better.
Update a single matrix $d$ \emph{in place}, overwriting $d^{(k-1)}$ with $d^{(k)}$ as we sweep $k$. The
worry is that when we read $d(i,k)$ and $d(k,j)$ we might already have overwritten them with their
$k$-th-layer values. But look at what happens to row $k$ and column $k$ at stage $k$:
\[
  d^{(k)}(i,k) = \min\big(d^{(k-1)}(i,k),\, d^{(k-1)}(i,k)+d^{(k-1)}(k,k)\big) = d^{(k-1)}(i,k),
\]
using $d^{(k-1)}(k,k)=0$ (true with no negative cycle). The same holds for $d^{(k)}(k,j)=d^{(k-1)}(k,j)$.
So the $k$-th row and column \emph{do not change} when door $k$ opens, and reading them post- or
pre-overwrite gives the same number. The in-place update is safe, and memory drops to $\Th(n^2)$.

\begin{algorithm}[H]
\caption{Floyd--Warshall (distances, predecessors, negative-cycle flag)}
\label{ch09:alg-fw}
\DontPrintSemicolon
\KwIn{Vertices $1..n$; weight $w$; $w(i,j)=+\infty$ if $(i,j)\notin E$.}
\KwOut{Matrix $d$ with $d(i,j)=\delta(i,j)$; predecessor matrix $\pi$.}
\For{$i \leftarrow 1$ \KwTo $n$}{
  \For{$j \leftarrow 1$ \KwTo $n$}{
    $d(i,j) \leftarrow w(i,j)$;\quad $d(i,i)\leftarrow 0$\;
    \lIf{$(i,j)\in E$ \textbf{and} $i\neq j$}{$\pi(i,j)\leftarrow i$} \lElse{$\pi(i,j)\leftarrow \text{nil}$}
  }
}
\For{$k \leftarrow 1$ \KwTo $n$}{
  \For{$i \leftarrow 1$ \KwTo $n$}{
    \For{$j \leftarrow 1$ \KwTo $n$}{
      \If{$d(i,k) + d(k,j) < d(i,j)$}{
        $d(i,j) \leftarrow d(i,k) + d(k,j)$\;
        $\pi(i,j) \leftarrow \pi(k,j)$\;
      }
    }
  }
}
\For{$i \leftarrow 1$ \KwTo $n$}{
  \lIf{$d(i,i) < 0$}{\textbf{report} ``negative cycle reachable through $i$''}
}
\Return $d,\pi$\;
\end{algorithm}

The predecessor update $\pi(i,j)\leftarrow\pi(k,j)$ records ``the last edge into $j$ on the best
$i$-to-$j$ path now enters $j$ the way the best $k$-to-$j$ path does.'' Following $\pi$ backward from
$j$ reconstructs the path, exactly as in single-source algorithms.

\subsection{A picture and a hand trace}
\label{ch09:fw-trace}

Take a five-vertex directed graph with one negative edge (but no negative cycle).

\begin{center}
\begin{tikzpicture}[>=Latex, node distance=24mm,
   v/.style={circle,draw=nodeblue,very thick,fill=blue!4,minimum size=8mm,font=\small}]
  \node[v] (1) {1};
  \node[v] (2) [right=of 1] {2};
  \node[v] (3) [right=of 2] {3};
  \node[v] (4) [below=18mm of 2] {4};
  \node[v] (5) [below=18mm of 3] {5};
  \draw[->,edgegray,thick] (1) -- node[above,font=\scriptsize]{3} (2);
  \draw[->,edgegray,thick] (2) -- node[above,font=\scriptsize]{4} (3);
  \draw[->,edgegray,thick] (1) -- node[left,font=\scriptsize]{8} (4);
  \draw[->,edgegray,thick] (4) -- node[below,font=\scriptsize]{$-2$} (2);
  \draw[->,edgegray,thick] (2) -- node[right,font=\scriptsize]{1} (5);
  \draw[->,edgegray,thick] (5) -- node[below,font=\scriptsize]{2} (3);
  \draw[->,edgegray,thick] (4) to[bend right=18] node[below,font=\scriptsize]{5} (5);
\end{tikzpicture}
\end{center}

Edges: $1\!\to\!2(3)$, $2\!\to\!3(4)$, $1\!\to\!4(8)$, $4\!\to\!2(-2)$, $2\!\to\!5(1)$,
$5\!\to\!3(2)$, $4\!\to\!5(5)$. The starting matrix $d^{(0)}$ (rows = from, columns = to, $\cdot$
means $+\infty$):

\[
d^{(0)}=
\begin{array}{c|ccccc}
 & 1 & 2 & 3 & 4 & 5\\\hline
1 & 0 & 3 & \cdot & 8 & \cdot\\
2 & \cdot & 0 & 4 & \cdot & 1\\
3 & \cdot & \cdot & 0 & \cdot & \cdot\\
4 & \cdot & -2 & \cdot & 0 & 5\\
5 & \cdot & \cdot & 2 & \cdot & 0
\end{array}
\]

\textbf{Open door $k=1$.} Vertex 1 has no incoming edges, so no path can route through it; nothing
changes. $d^{(1)}=d^{(0)}$.

\textbf{Open door $k=2$.} Now paths may detour through 2. Row 4 reaches 2 cheaply via the $-2$ edge:
$d(4,3)\leftarrow d(4,2)+d(2,3)=-2+4=2$, and $d(4,5)\leftarrow\min(5,\,d(4,2)+d(2,5))=\min(5,-2+1)=-1$.
Row 1 already reached 2 directly, and through-2 paths give $d(1,3)\leftarrow d(1,2)+d(2,3)=3+4=7$ and
$d(1,5)\leftarrow 3+1=4$.
\[
d^{(2)}=
\begin{array}{c|ccccc}
 & 1 & 2 & 3 & 4 & 5\\\hline
1 & 0 & 3 & 7 & 8 & 4\\
2 & \cdot & 0 & 4 & \cdot & 1\\
3 & \cdot & \cdot & 0 & \cdot & \cdot\\
4 & \cdot & -2 & 2 & 0 & -1\\
5 & \cdot & \cdot & 2 & \cdot & 0
\end{array}
\]

\textbf{Open door $k=3$.} Vertex 3 is a sink (no outgoing edges), so it is never useful as an
intermediate. $d^{(3)}=d^{(2)}$.

\textbf{Open door $k=4$.} Who can reach 4? Only vertex 1 ($d(1,4)=8$). Routing 1 through 4 gives
$d(1,2)\leftarrow\min(3,\,8+(-2))=3$, $d(1,3)\leftarrow\min(7,\,8+2)=7$,
$d(1,5)\leftarrow\min(4,\,8+(-1))=4$ -- no improvements. Door 4 helps nobody here.

\textbf{Open door $k=5$.} Paths may detour through 5 to reach 3 cheaply ($5\to3$ costs 2):
$d(2,3)\leftarrow\min(4,\,d(2,5)+d(5,3))=\min(4,\,1+2)=3$,
$d(4,3)\leftarrow\min(2,\,d(4,5)+d(5,3))=\min(2,\,-1+2)=1$, and
$d(1,3)\leftarrow\min(7,\,d(1,5)+d(5,3))=\min(7,\,4+2)=6$.
\[
d^{(5)}=
\begin{array}{c|ccccc}
 & 1 & 2 & 3 & 4 & 5\\\hline
1 & 0 & 3 & 6 & 8 & 4\\
2 & \cdot & 0 & 3 & \cdot & 1\\
3 & \cdot & \cdot & 0 & \cdot & \cdot\\
4 & \cdot & -2 & 1 & 0 & -1\\
5 & \cdot & \cdot & 2 & \cdot & 0
\end{array}
\]
This is the final $\delta$ matrix. Sanity-check $\delta(1,3)=6$: the path $1\to2\to5\to3$ costs
$3+1+2=6$; the direct-looking $1\to2\to3$ costs $3+4=7$. The cheaper one wins, and it is exactly the
path the $k=5$ door opened. The negative edge $4\to2$ never broke anything -- Floyd--Warshall handles
it without complaint.

\begin{labbox}
Re-run the trace, but add the edge $3\to1$ with weight $-10$. After $k$ reaches the vertices on the
cycle $1\to2\to3\to1$ (cost $3+4-10=-3$), watch $d(1,1)$ drop below $0$. That negative diagonal entry
is the algorithm waving a red flag: a negative cycle is reachable from and back to vertex 1, so the
distances among those vertices are $-\infty$. Which diagonal entries go negative, and at which value
of $k$?
\end{labbox}

\subsection{Correctness}

\begin{theorem}[Floyd--Warshall correctness]
\label{ch09:thm-fw}
If $\Gph$ has no negative-weight cycle, then after Algorithm~\ref{ch09:alg-fw} terminates
$d(i,j)=\delta(i,j)$ for every pair $(i,j)$.
\end{theorem}

\begin{proof}
We prove by induction on $k$ that after the outer loop has processed door $k$, the matrix holds
$d^{(k)}(i,j)$: the shortest-path cost from $i$ to $j$ using intermediates only in $\set{1,\dots,k}$.

\emph{Base case} $k=0$: the initialization sets $d(i,j)=w(i,j)$ (and $0$ on the diagonal), which is
$d^{(0)}$ by definition -- single edges are the only intermediate-free paths.

\emph{Inductive step.} Assume the matrix holds $d^{(k-1)}$ before door $k$ is processed. Fix $(i,j)$
and consider an optimal path $P^\star$ from $i$ to $j$ with intermediates in $\set{1,\dots,k}$; since
there is no negative cycle we may take $P^\star$ to be simple. If $k\notin P^\star$ then $P^\star$
uses only $\set{1,\dots,k-1}$ and has cost $d^{(k-1)}(i,j)$. If $k\in P^\star$, splitting at $k$ gives
two subpaths with interiors in $\set{1,\dots,k-1}$, each itself optimal among such paths (an exchange
argument: a cheaper subpath would yield a cheaper $P^\star$). Their costs are $d^{(k-1)}(i,k)$ and
$d^{(k-1)}(k,j)$. Hence the true optimum equals
$\min(d^{(k-1)}(i,j),\,d^{(k-1)}(i,k)+d^{(k-1)}(k,j))$, which is exactly what the inner update
(equation~\ref{ch09:fw-recurrence}) computes. By the in-place argument above, reading $d(i,k)$ and
$d(k,j)$ returns their $d^{(k-1)}$ values even though we update in place. So after door $k$ the matrix
holds $d^{(k)}$. At $k=n$ this is $\delta$.
\end{proof}

\subsection{Negative-cycle detection}
\label{ch09:fw-negcycle}

\begin{theorem}
\label{ch09:thm-fw-negcycle}
Run Algorithm~\ref{ch09:alg-fw} on \emph{any} weighted digraph (initializing the diagonal to $0$).
After the triple loop, $\Gph$ contains a negative-weight cycle if and only if $d(i,i)<0$ for some $i$.
\end{theorem}

\begin{proof}
($\Leftarrow$) If $d(i,i)<0$ then the recurrence found a closed walk from $i$ back to $i$ of negative
total weight; a closed walk of negative weight contains a negative cycle.

($\Rightarrow$) Let $C=v_0\to v_1\to\cdots\to v_t=v_0$ be a negative cycle and let $v_p$ be its
highest-numbered vertex. Rotate $C$ so it starts and ends at $v_p$. Every other vertex of $C$ has index
$<v_p$, so once door $v_p$ is processed both subpaths $v_p\rightsquigarrow$ and $\rightsquigarrow v_p$
have interiors below $v_p$, and the update sets $d(v_p,v_p)\leftarrow\min(0,\,\text{negative total})<0$.
Once a diagonal goes negative the $\min$ keeps it negative. So some diagonal entry is negative at the
end.
\end{proof}

\begin{trapbox}
With negative cycles present, the \emph{off}-diagonal entries are not meaningful: a value like
$d(i,j)=7$ might really be $-\infty$ because a negative cycle sits between $i$ and $j$. The clean fix is
a post-pass: for every $k$ with $d(k,k)<0$, set $d(i,j)\leftarrow-\infty$ whenever $d(i,k)<\infty$ and
$d(k,j)<\infty$ (you can reach $k$ and leave it). Only after this pass should you trust the off-diagonal
numbers.
\end{trapbox}

\begin{complexitybox}
Floyd--Warshall runs in $\Th(n^3)$ time and $\Th(n^2)$ space, independent of $m$. Three nested loops,
one comparison and possible addition each. No priority queue, no recursion, and the inner loop is a
tight scan over contiguous rows -- it vectorizes and caches beautifully, which is why it routinely beats
$n$ Dijkstras on dense graphs even though the asymptotics tie.
\end{complexitybox}

\section{Transitive closure: reachability with the same loop}
\label{ch09:closure}

Strip the weights away and ask a yes/no question: is there \emph{any} path from $i$ to $j$? The answer
is the transitive closure $\Gph^{*}$, the reachability matrix $R$ with $R(i,j)=1$ iff $j$ is reachable
from $i$. This is Warshall's original 1962 algorithm, and it is Floyd--Warshall with $(\min,+)$ replaced
by $(\textsc{or},\textsc{and})$:
\[
  R^{(k)}(i,j) = R^{(k-1)}(i,j)\;\;\textsc{or}\;\;\big(R^{(k-1)}(i,k)\;\textsc{and}\;R^{(k-1)}(k,j)\big).
\]
``$i$ reaches $j$ already, or $i$ reaches $k$ and $k$ reaches $j$.'' Same door-by-door logic.

\begin{algorithm}[H]
\caption{Warshall transitive closure}
\label{ch09:alg-warshall}
\DontPrintSemicolon
$R(i,j)\leftarrow [\,(i,j)\in E\,]$ for all $i,j$;\quad $R(i,i)\leftarrow 1$\;
\For{$k\leftarrow 1$ \KwTo $n$}{
  \For{$i\leftarrow 1$ \KwTo $n$}{
    \lIf{$R(i,k)$}{\For{$j\leftarrow 1$ \KwTo $n$}{$R(i,j)\leftarrow R(i,j)\;\textsc{or}\;R(k,j)$}}
  }
}
\Return $R$\;
\end{algorithm}

The guard ``\textbf{if} $R(i,k)$'' hoists the test for $k$ out of the innermost loop -- when $i$ cannot
reach $k$, the whole $j$-scan is skipped.

\begin{whybox}
Why does the same triple loop solve both problems? Because both are instances of the same algebraic
template: a \emph{closed semiring}. Shortest paths use $(\R\cup\set{\infty},\min,+)$ with identities
$\infty$ and $0$; reachability uses $(\set{0,1},\textsc{or},\textsc{and})$ with identities $0$ and $1$.
Floyd--Warshall is really the Gauss--Jordan / Kleene-star computation over a semiring, and any semiring
with the right absorption properties drops straight in -- bottleneck paths use $(\min,\max)$, path
counting uses $(+,\times)$, and regular-expression construction (Kleene's theorem) uses $(\cup,\cdot)$
over languages.
\end{whybox}

\begin{applicationbox}
Bit-packing makes closure fast in practice. Store each row of $R$ as a machine-word bitset. The inner
``$R(i,\cdot)\mathrel{|}=R(k,\cdot)$'' becomes a word-parallel OR over $\ceil{n/w}$ words, cutting the
running time to $\Th(n^3/w)$ for word size $w=64$. This is the standard trick for reachability in
compilers (which definitions reach a program point) and for the meet-in-the-middle phase of many graph
queries.
\end{applicationbox}

\section{The matrix-multiplication view and the min-plus product}
\label{ch09:minplus}

There is a second, more algebraic, dynamic program for APSP, and it exposes a deep structure. Define
\[
  L^{(\ell)}(i,j) = \text{weight of the shortest walk from } i \text{ to } j \text{ using at most } \ell \text{ edges.}
\]
A walk with at most $\ell$ edges is a walk with at most $\ell-1$ edges followed by one more edge:
\[
  L^{(\ell)}(i,j) = \min_{1\le k\le n}\;\Big(L^{(\ell-1)}(i,k) + w(k,j)\Big),
  \qquad L^{(1)}(i,j)=w(i,j).
\]
Stare at that formula. It is matrix multiplication with $+$ replaced by $\min$ and $\times$ replaced by
$+$. This is the \textbf{min-plus product} (also called the \emph{tropical} or \emph{distance} product).
Writing $\otimes$ for it,
\[
  (A\otimes B)(i,j) = \min_{k}\big(A(i,k)+B(k,j)\big),
  \qquad L^{(\ell)} = L^{(\ell-1)}\otimes W = W^{\otimes \ell}.
\]

\begin{intuitionbox}
Ordinary matrix powering counts walks: $(\Adjm^\ell)(i,j)$ is the number of length-$\ell$ walks from $i$
to $j$, because $\sum_k$ adds up ways and $\times$ combines independent choices. Swap the operations to
$(\min,+)$ and ``count walks'' becomes ``find the lightest walk'' -- summation over ways turns into
minimization over alternatives, and multiplication of counts turns into addition of lengths. Same
skeleton, tropical flesh.
\end{intuitionbox}

In a graph with no negative cycle, every shortest path is simple and uses at most $n-1$ edges, so
$L^{(n-1)}=\delta$. Computing $W^{\otimes(n-1)}$ by repeated squaring needs only
$\ceil{\log_2(n-1)}$ min-plus products:
\[
  W^{\otimes 1},\; W^{\otimes 2},\; W^{\otimes 4},\;\dots,\; W^{\otimes 2^{\lceil\log_2(n-1)\rceil}}.
\]
Each $\otimes$ is the obvious triple loop at $\Th(n^3)$, so this APSP route costs $\Th(n^3\log n)$ -- a
$\log n$ factor \emph{worse} than Floyd--Warshall. Why bother?

\begin{aestheticbox}
Two payoffs. First, the min-plus formulation reveals APSP as a linear-algebra problem in disguise, which
is where the asymptotically fastest known algorithms live. If one could compute a min-plus product in
$\Oh(n^{3-\eps})$ time, APSP would fall below cubic; the long-standing \emph{APSP conjecture} holds that
no truly subcubic ($\Oh(n^{3-\eps})$) combinatorial algorithm exists, and a web of problems (negative
triangle, min-plus product, APSP) are known to be subcubic-equivalent. Williams' bound of
$n^3/2^{\Th(\sqrt{\log n})}$ shaves off subpolynomial factors but not a true power. Second, over a
\emph{small integer} weight range the min-plus product can be embedded into a few real matrix
multiplications and accelerated with Strassen-style $\Oh(n^\omega)$ algebra (Seidel, Zwick), giving
genuinely subcubic APSP on unweighted or small-weight graphs.
\end{aestheticbox}

\begin{historybox}
Min-plus algebra is the \emph{tropical semiring}, named (half-jokingly) after the Brazilian mathematician
Imre Simon. Seidel's 1995 algorithm computes unweighted-undirected APSP in $\Oh(n^\omega\log n)$ by
squaring the adjacency matrix and recovering distances via parities; Zwick (2002) pushed weighted
directed APSP to about $\Oh(n^{2.58})$ for small integer weights using rectangular matrix multiplication.
The combinatorial frontier, $n^3/2^{\Th(\sqrt{\log n})}$ (Williams 2014), remains the best general
bound -- evidence for the conjectured cubic barrier.
\end{historybox}

\begin{connectionbox}
The min-plus / matrix-power lens is your bridge to the linear-algebra companion. Counting walks with
$\Adjm^\ell$ links to the spectrum of the adjacency matrix (the spectral graph theory chapter); the
resolvent $(I-x\Adjm)^{-1}=\sum_\ell x^\ell \Adjm^\ell$ is a walk generating function. Replace the field
$(\R,+,\times)$ by the semiring $(\R\cup\set{\infty},\min,+)$ and matrix \emph{inversion} becomes the
Kleene star $W^*=\bigoplus_{\ell\ge 0}W^{\otimes\ell}$ -- which is precisely the all-pairs distance
matrix. Floyd--Warshall is Gaussian elimination over this semiring. The Laplacian and electrical-flow
chapters reuse this ``same skeleton, different algebra'' idea when they replace combinatorial distance by
effective resistance.
\end{connectionbox}

\section{Johnson's algorithm: reweighting for sparse graphs}
\label{ch09:johnson}

Floyd--Warshall is $\Th(n^3)$ no matter how few edges the graph has. On a sparse graph ($m=\Oh(n)$ or
so) we would much rather pay $\Oh(nm+n^2\log n)$ -- $n$ runs of Dijkstra. The obstacle is negative edges:
Dijkstra needs nonnegative weights. Johnson's algorithm removes the obstacle by \emph{reweighting}.

\subsection{The potential trick}

\begin{intuitionbox}
Assign every vertex $v$ a number $h(v)$ -- a \emph{height} or \emph{potential}. Reweight each edge by how
much it goes uphill:
\[
  \hat{w}(u,v) = w(u,v) + h(u) - h(v).
\]
Now walk any path $u_0\to u_1\to\cdots\to u_t$. The heights telescope:
\[
  \hat{w}(\text{path}) = \sum_{i}\big(w(u_{i-1},u_i)+h(u_{i-1})-h(u_i)\big)
   = w(\text{path}) + h(u_0) - h(u_t).
\]
Every path between a fixed pair $(u_0,u_t)$ shifts by the \emph{same} constant $h(u_0)-h(u_t)$. So
reweighting preserves the \emph{ordering} of paths by cost -- the shortest path stays the shortest. We
have not changed the problem, only the bookkeeping.
\end{intuitionbox}

We want heights that make every reweighted edge nonnegative:
\[
  \hat{w}(u,v) = w(u,v)+h(u)-h(v) \ge 0
  \quad\Longleftrightarrow\quad
  h(v) - h(u) \le w(u,v) \quad\text{for every edge } (u,v).
\]
These are exactly the constraints satisfied by shortest-path \emph{distances} from any source: if
$h(v)=\delta(s,v)$ then the triangle inequality $\delta(s,v)\le\delta(s,u)+w(u,v)$ gives precisely
$h(v)-h(u)\le w(u,v)$. So a single-source shortest-path tree hands us a feasible potential -- as long as
we can compute it, which we can with Bellman--Ford, the one solver that tolerates negative edges.

\begin{trapbox}
There is a catch: a source $s$ might not reach every vertex, leaving $\delta(s,v)=\infty$ and a useless
potential. Johnson's fix is a \emph{virtual super-source} $q$ with a zero-weight edge $q\to v$ to every
vertex. Now $q$ reaches everyone, $h(v)=\delta(q,v)$ is finite for all $v$, and the new edges create no
negative cycle (they point only outward from $q$, which has no incoming edges). If Bellman--Ford from $q$
detects a negative cycle, we abort -- APSP is ill-posed.
\end{trapbox}

\subsection{The algorithm}

\begin{algorithm}[H]
\caption{Johnson's all-pairs shortest paths}
\label{ch09:alg-johnson}
\DontPrintSemicolon
\KwIn{Digraph $\Gph=(V,E)$, weights $w$ (possibly negative).}
\KwOut{$\delta(u,v)$ for all $u,v$, or report a negative cycle.}
\textbf{Step 1 (super-source).} Build $\Gph'$ by adding $q$ and edges $q\to v$ of weight $0$ for all $v\in V$\;
\textbf{Step 2 (potentials).} Run Bellman--Ford from $q$ on $\Gph'$\;
\lIf{a negative cycle is found}{\textbf{report} and \textbf{stop}}
$h(v)\leftarrow \delta_{\Gph'}(q,v)$ for every $v\in V$\;
\textbf{Step 3 (reweight).} \For{each edge $(u,v)\in E$}{$\hat{w}(u,v)\leftarrow w(u,v)+h(u)-h(v)$ \tcp*{$\ge 0$}}
\textbf{Step 4 (Dijkstra from each source).}\;
\For{each $u\in V$}{
  run Dijkstra from $u$ with weights $\hat{w}$ to get $\hat{\delta}(u,v)$ for all $v$\;
  \For{each $v\in V$}{$\delta(u,v)\leftarrow \hat{\delta}(u,v) - h(u) + h(v)$ \tcp*{undo the shift}}
}
\Return $\delta$\;
\end{algorithm}

The undo step inverts the telescoping: since $\hat\delta(u,v)=\delta(u,v)+h(u)-h(v)$, we recover
$\delta(u,v)=\hat\delta(u,v)-h(u)+h(v)$.

\subsection{Worked example: same graph, reweighted}
\label{ch09:johnson-trace}

Reuse the five-vertex graph from Section~\ref{ch09:fw-trace}, whose only negative edge is $4\to2(-2)$.
Add super-source $q$ with zero edges to all five vertices (Figure~\ref{ch09:fig-supersource}) and run
Bellman--Ford from $q$. Each $q\to v$ edge has weight $0$, so $h(v)$ is the lightest way to reach $v$
from anywhere (through $q$). Relaxation propagates through the negative edge, so we follow the chains
carefully:

\begin{figure}[H]
\centering
\begin{tikzpicture}[>=Latex, node distance=20mm,
   v/.style={circle,draw=nodeblue,very thick,fill=blue!4,minimum size=8mm,font=\small},
   q/.style={circle,draw=quotientorange,very thick,fill=orange!8,minimum size=8mm,font=\small}]
  \node[v] (1) {1};
  \node[v] (2) [right=of 1] {2};
  \node[v] (3) [right=of 2] {3};
  \node[v] (4) [below=16mm of 2] {4};
  \node[v] (5) [below=16mm of 3] {5};
  \node[q] (q) [left=20mm of 4] {$q$};
  \draw[->,edgegray,thick] (1) -- node[above,font=\scriptsize]{3} (2);
  \draw[->,edgegray,thick] (2) -- node[above,font=\scriptsize]{4} (3);
  \draw[->,edgegray,thick] (1) -- node[left,font=\scriptsize]{8} (4);
  \draw[->,edgegray,thick] (4) -- node[below,font=\scriptsize]{$-2$} (2);
  \draw[->,edgegray,thick] (2) -- node[right,font=\scriptsize]{1} (5);
  \draw[->,edgegray,thick] (5) -- node[below,font=\scriptsize]{2} (3);
  \draw[->,edgegray,thick] (4) to[bend right=18] node[below,font=\scriptsize]{5} (5);
  \draw[->,quotientorange,thick,dashed] (q) to[bend left=10] node[above,font=\scriptsize]{0} (1);
  \draw[->,quotientorange,thick,dashed] (q) to node[above,font=\scriptsize,pos=0.85]{0} (2);
  \draw[->,quotientorange,thick,dashed] (q) to[bend left=14] node[above,font=\scriptsize,pos=0.9]{0} (3);
  \draw[->,quotientorange,thick,dashed] (q) -- node[below,font=\scriptsize]{0} (4);
  \draw[->,quotientorange,thick,dashed] (q) to[bend right=10] node[below,font=\scriptsize,pos=0.85]{0} (5);
\end{tikzpicture}
\caption{The worked graph with Johnson's super-source $q$ (dashed orange edges of weight $0$ to every
vertex). $q$ has no incoming edges, so it adds no cycle; Bellman--Ford from $q$ yields the
potentials $h(v)=\delta(q,v)$.}
\label{ch09:fig-supersource}
\end{figure}

\begin{itemize}
  \item $h(1)=0$ (direct $q\to1$; nothing reaches $1$).
  \item $h(2)=-2$ via $q\to4\to2$, beating the direct $q\to2=0$.
  \item $h(5)=-1$ via $q\to4\to2\to5 = 0-2+1$.
  \item $h(3)=1$ via $q\to4\to2\to5\to3 = 0-2+1+2$.
  \item $h(4)=0$ (direct $q\to4$; nothing else reaches $4$).
\end{itemize}
\[
\begin{array}{c|ccccc}
 v & 1 & 2 & 3 & 4 & 5\\\hline
 h(v) & 0 & -2 & 1 & 0 & -1
\end{array}
\]
Now reweight every original edge with $\hat{w}(u,v)=w(u,v)+h(u)-h(v)$:
\[
\begin{array}{l|c|c}
\text{edge } (u,v) & w & \hat{w}=w+h(u)-h(v)\\\hline
1\to2 & 3 & 3+0-(-2)=5\\
2\to3 & 4 & 4+(-2)-1=1\\
1\to4 & 8 & 8+0-0=8\\
4\to2 & -2 & -2+0-(-2)=0\\
2\to5 & 1 & 1+(-2)-(-1)=0\\
5\to3 & 2 & 2+(-1)-1=0\\
4\to5 & 5 & 5+0-(-1)=6
\end{array}
\]
Every reweighted edge is nonnegative -- exactly as Theorem~\ref{ch09:thm-johnson} promises -- and the
negative edge $4\to2$ became a harmless $0$.

Run Dijkstra from source $1$ on $\hat w$. It settles $\hat\delta(1,2)=5$, $\hat\delta(1,4)=8$, and
$\hat\delta(1,5)=\min(5+0,\,8+6)=5$ (route $1\to2\to5$). For $1\to3$ the contenders are
$1\to2\to3$ at $5+1=6$ and $1\to2\to5\to3$ at $5+0+0=5$, so $\hat\delta(1,3)=5$. Undo the shift for
the $1\to3$ answer:
\[
  \delta(1,3) = \hat\delta(1,3) - h(1) + h(3) = 5 - 0 + 1 = 6,
\]
matching the $\delta(1,3)=6$ from Floyd--Warshall. Likewise
$\delta(1,5)=\hat\delta(1,5)-h(1)+h(5)=5-0+(-1)=4$ and $\delta(1,2)=5-0+(-2)=3$. Every number agrees.

\begin{labbox}
Finish the trace: run Dijkstra from source $4$ on $\hat w$ (starting across the $4\to2(0)$ edge) and undo
the shift with $h(4)=0$. You should recover row $4$ of the Floyd--Warshall matrix:
$\delta(4,2)=-2$, $\delta(4,3)=1$, $\delta(4,5)=-1$. Notice the reweighted Dijkstra distances are all
$\ge 0$ even though the true distances are negative -- the potentials carry the negativity.
\end{labbox}

\subsection{Correctness}

\begin{theorem}[Johnson reweighting]
\label{ch09:thm-johnson}
Let $h(v)=\delta_{\Gph'}(q,v)$ be finite for all $v$ (guaranteed by the super-source when no negative
cycle exists). Then:
\begin{enumerate}
  \item[(a)] $\hat{w}(u,v)=w(u,v)+h(u)-h(v)\ge 0$ for every edge $(u,v)\in E$;
  \item[(b)] for every pair $(u,v)$, a path is shortest under $w$ iff it is shortest under $\hat w$, and
        $\delta(u,v)=\hat\delta(u,v)-h(u)+h(v)$.
\end{enumerate}
\end{theorem}

\begin{proof}
(a) Since $h(v)=\delta_{\Gph'}(q,v)$ is a shortest-path distance, it satisfies the triangle inequality
along the edge $(u,v)$: $\delta_{\Gph'}(q,v)\le\delta_{\Gph'}(q,u)+w(u,v)$, i.e. $h(v)\le h(u)+w(u,v)$.
Rearranged, $w(u,v)+h(u)-h(v)\ge 0$, which is $\hat w(u,v)\ge 0$.

(b) Take any path $P=(u=x_0,x_1,\dots,x_t=v)$. Summing $\hat w$ along $P$ telescopes:
\[
  \hat w(P)=\sum_{i=1}^{t}\big(w(x_{i-1},x_i)+h(x_{i-1})-h(x_i)\big)=w(P)+h(u)-h(v).
\]
The term $h(u)-h(v)$ depends only on the endpoints, not on the path. Hence for fixed $(u,v)$ the map
$P\mapsto\hat w(P)$ is $w(P)$ plus a constant, so it preserves the argmin: the same path minimizes both.
Taking minima, $\hat\delta(u,v)=\delta(u,v)+h(u)-h(v)$, and rearranging gives the recovery formula.
Because every $\hat w\ge 0$, Dijkstra correctly computes $\hat\delta$.
\end{proof}

\begin{whybox}
The potential $h$ is a \emph{certificate}, not magic. Feasibility ``$h(v)-h(u)\le w(u,v)$ for all edges''
is the linear-programming dual of shortest paths; a feasible potential exists \emph{iff} there is no
negative cycle (Bellman--Ford either finds one or proves none exists). So Johnson's algorithm secretly
solves the shortest-path LP's dual once, then exploits that dual solution $n$ times. The same reweighting
reappears in the successive-shortest-paths method for min-cost flow, where node potentials keep residual
edges nonnegative so each augmentation can use Dijkstra.
\end{whybox}

\begin{complexitybox}
Johnson's algorithm: Bellman--Ford once at $\Oh(nm)$, reweighting at $\Oh(m)$, then $n$ Dijkstra runs.
With a binary heap each Dijkstra is $\Oh((n+m)\log n)$, total $\Oh(nm\log n)$; with a Fibonacci heap each
is $\Oh(m+n\log n)$, total $\Oh(nm+n^2\log n)$. On sparse graphs $m=\Oh(n)$ this is $\Oh(n^2\log n)$ --
asymptotically beating Floyd--Warshall's $\Th(n^3)$. On dense graphs $m=\Th(n^2)$ both are $\Th(n^3)$
(Fibonacci) and Floyd--Warshall wins on constants.
\end{complexitybox}

\section{Choosing an algorithm}
\label{ch09:choosing}

\begin{center}
\begin{tabular}{@{}llll@{}}
\toprule
Situation & Algorithm & Time & Notes\\
\midrule
Dense, any weights & Floyd--Warshall & $\Th(n^3)$ & simplest; cache-friendly\\
Sparse, nonnegative & $n\times$ Dijkstra & $\Oh(nm+n^2\log n)$ & no reweighting needed\\
Sparse, negative edges & Johnson & $\Oh(nm+n^2\log n)$ & reweight, then $n$ Dijkstras\\
Reachability only & Warshall (bitset) & $\Th(n^3/w)$ & word-parallel OR\\
Small integer weights & Seidel / Zwick & $\Oh(n^\omega\,\poly\log)$ & via fast matmul\\
\bottomrule
\end{tabular}
\end{center}

\begin{intuitionbox}
A one-line decision rule. If the graph is dense or you just want the simplest correct code, write
Floyd--Warshall. If it is sparse, run Dijkstra from every vertex -- and if any edge is negative, prefix
that with one Johnson reweighting pass. Everything else is a refinement.
\end{intuitionbox}

\section{Path reconstruction and a few extensions}
\label{ch09:extensions}

\paragraph{Recovering the actual paths.} Both algorithms can return paths, not just costs. In
Floyd--Warshall the predecessor matrix $\pi$ (Algorithm~\ref{ch09:alg-fw}) lets you walk backward from
$j$: repeatedly replace $j$ by $\pi(i,j)$ until you reach $i$. In Johnson's algorithm each Dijkstra
already produces a shortest-path tree; reweighting changes edge costs but not path \emph{structure}
(Theorem~\ref{ch09:thm-johnson}(b)), so the tree edges are the genuine shortest paths.

\paragraph{Bottleneck / widest paths.} Replace $(\min,+)$ by $(\min,\max)$ in the Floyd--Warshall
update -- $d(i,j)\leftarrow\min(d(i,j),\,\max(d(i,k),d(k,j)))$ -- and you compute the minimax path between
every pair, the path whose heaviest edge is as light as possible. This is the all-pairs minimum-bottleneck
problem (and with $(\max,\min)$, maximum-capacity), useful in network routing.

\paragraph{Counting shortest paths.} Carry a count matrix $c$ alongside $d$. When the update strictly
improves $d(i,j)$, set $c(i,j)\leftarrow c(i,k)\cdot c(k,j)$; when it \emph{ties} the current best, add
$c(i,j)\mathrel{+}= c(i,k)\cdot c(k,j)$. This computes the number of distinct shortest paths between every
pair, a quantity central to betweenness-centrality computations on networks.

\begin{applicationbox}
APSP distance matrices feed graph-analytic measures directly. The \emph{Wiener index}
$\sum_{i<j}\delta(i,j)$ measures molecular branching in chemistry. \emph{Closeness centrality}
$1/\sum_j\delta(i,j)$ ranks how central a vertex is; \emph{betweenness} counts shortest paths through each
vertex (the counting extension above). The graph \emph{diameter} $\max_{i,j}\delta(i,j)$ and \emph{radius}
fall out of one scan of the matrix. Whenever a pipeline needs ``distance between everything,'' an APSP
solver is the engine.
\end{applicationbox}

\section{Exercises}
\label{ch09:exercises}

\begin{exercise}
Implement Floyd--Warshall with the predecessor matrix $\pi$ of Algorithm~\ref{ch09:alg-fw} and a path
reconstruction routine. Test it on the five-vertex graph of Section~\ref{ch09:fw-trace} and confirm it
reports the path $1\to2\to5\to3$ for the pair $(1,3)$.
\end{exercise}

\begin{exercise}
Prove that the in-place update is safe by showing rigorously that row $k$ and column $k$ of $d$ are
unchanged during the processing of door $k$, and explain exactly where the assumption $d^{(k-1)}(k,k)=0$
is used. What goes wrong if a negative cycle through $k$ exists?
\end{exercise}

\begin{exercise}
Add the edge $3\to1$ with weight $-10$ to the worked graph, creating the cycle $1\to2\to3\to1$ of weight
$-3$. Trace Floyd--Warshall and identify the first value of $k$ at which some diagonal entry $d(i,i)$
becomes negative. Then implement the $-\infty$ post-pass from the trap box in
Section~\ref{ch09:fw-negcycle} and report which off-diagonal entries it corrects.
\end{exercise}

\begin{exercise}
(Min-plus by squaring.) Implement APSP via repeated squaring of the min-plus product
(Section~\ref{ch09:minplus}). Verify it agrees with Floyd--Warshall on random graphs, and confirm
empirically that it performs $\ceil{\log_2(n-1)}$ products. Why does ordinary squaring fail to detect
negative cycles the way Floyd--Warshall's diagonal does?
\end{exercise}

\begin{exercise}
Prove that the min-plus product is associative, so that $W^{\otimes a}\otimes W^{\otimes b}=W^{\otimes(a+b)}$
and repeated squaring is valid. Where in the proof do you use that $(\R\cup\set{\infty},\min,+)$ is a
semiring (in particular, that $+$ distributes over $\min$)?
\end{exercise}

\begin{exercise}
Show that a feasible potential $h$ with $h(v)-h(u)\le w(u,v)$ for all edges exists if and only if $\Gph$
has no negative cycle. (Hint: one direction is Johnson's super-source construction; for the other, sum
the inequalities around a cycle and watch the potentials cancel.)
\end{exercise}

\begin{exercise}
In Johnson's algorithm, suppose you skip the super-source and instead pick an arbitrary vertex $s$ as the
Bellman--Ford source. Construct a small graph where this produces an infinite potential for some vertex
and breaks the reweighting. Then argue that the super-source never introduces a negative cycle.
\end{exercise}

\begin{exercise}
(Bottleneck APSP.) Modify Floyd--Warshall to compute, for every pair, the minimum possible value of the
\emph{maximum} edge weight on a path (the minimax path). Prove your recurrence correct by an exchange
argument analogous to Theorem~\ref{ch09:thm-fw}, and state the semiring you are working over.
\end{exercise}

\begin{exercise}
(Counting shortest paths.) Augment Floyd--Warshall with the count matrix $c$ described in
Section~\ref{ch09:extensions}. Prove that on a graph with nonnegative weights and no zero-weight cycles
the final $c(i,j)$ equals the number of distinct shortest paths from $i$ to $j$. Explain why zero-weight
cycles can make this count infinite.
\end{exercise}

\begin{exercise}
(Transitive closure as a special case.) Show that running Floyd--Warshall with edge weights $0$ on
present edges and $+\infty$ on absent ones, then thresholding $d(i,j)<\infty$, recovers the transitive
closure. Compare its constant factor with the bit-packed Warshall closure of
Algorithm~\ref{ch09:alg-warshall} and explain when the bitset version wins.
\end{exercise}

\begin{exercise}
(Incremental APSP.) After computing the full distance matrix, a single edge weight $w(a,b)$ is
\emph{decreased}. Give an $\Oh(n^2)$ algorithm to update all $n^2$ distances, and prove it correct.
(Hint: the only paths that can improve route through $(a,b)$; relax
$d(i,j)\leftarrow\min(d(i,j),\,d(i,a)+w(a,b)+d(b,j))$.) Why does an edge \emph{increase} not admit the
same easy update?
\end{exercise}
